\section*{Chapitre \ref{ch:s1:seasonality-and-calendar-effects} --- Saisonnalité et effets calendaires}

\begin{solution}{exo:s1:seasonality-and-calendar-effects:1}
$100 \times 0.05 = 5$~; avec 75 règles nulles, $75 \times 0.05 = 3.75$. L'expérience en a trouvé 8, ce qui reste dans le bruit des règles corrélées (un faux jour de la semaine
tend à entraîner avec lui de faux jours et de fausses quinzaines).
\end{solution}

\begin{solution}{exo:s1:seasonality-and-calendar-effects:2}
$\Phi^{-1}(1 - 0.05/200) = 3.48$.
\end{solution}

\begin{solution}{exo:s1:seasonality-and-calendar-effects:3}
Personne ne peut acheter du gaz de janvier en juillet et le garder au comptant~; les instruments négociables sont les contrats à terme, dont les cours pour une livraison en hiver incluent déjà la
prime hivernale attendue. Une position sur contrats à terme ne rapporte que l'écart entre la saison réalisée et celle déjà intégrée au cours, auquel s'ajoute toute prime de risque.
\end{solution}

\begin{solution}{exo:s1:seasonality-and-calendar-effects:4}
La volatilité quotidienne vaut $0.147/\sqrt{252} = 0.93\%$. Pour l'\omterm{def:s1:seasonality-and-calendar-effects:tom}{effet de fin de mois}, l'erreur type vaut $0.93\% \times \sqrt{1/1\,440 + 1/6\,120} = 2.71$ points
de base, ce qui donne en moyenne $t = 2.95$ pour 8 points de base (l'échantillon tiré a donné 4.4). Pour les jours fériés, elle vaut $0.93~\% \times \sqrt{1/270 + 1/7\,290} = 5.74$
points de base et $t = 2.61$ (l'échantillon a donné 1.5).
\end{solution}

\begin{solution}{exo:s1:seasonality-and-calendar-effects:5}
La statistique de Student croît avec la racine carrée de la taille de l'échantillon~: soit $30 \times (3.48/2.61)^2 = 53$ années.
\end{solution}

\begin{solution}{exo:s1:seasonality-and-calendar-effects:6}
Il contrôle la part attendue de fausses découvertes parmi celles faites, et non la probabilité d'en faire au moins une~; avec de nombreux vrais effets, il en conserve davantage et
tolère quelques fausses découvertes. Cela est acceptable lorsque les découvertes alimentent des tests ultérieurs (un pipeline de recherche), et non lorsque chacune est négociée dès qu'elle est trouvée.
\end{solution}

\begin{solution}{exo:s1:seasonality-and-calendar-effects:7}
Avec la graine~8, la meilleure règle de fin de mois a $t = 2.25$ et la règle de veille de jour férié $t = 2.99$~: sans correction, 8 règles de fin de mois, la règle de veille de jour férié
et 4 règles sans aucun effet réel sont retenues~; toute correction n'en retient aucune. Les effets injectés n'ont pas changé~; c'est l'échantillon qui a changé. Ce qui doit rester, c'est la
régularité des corrections (peu ou pas de fausses découvertes)~; ce qui change, c'est lesquels des vrais effets ont assez de chance pour passer.
\end{solution}

\begin{solution}{exo:s1:seasonality-and-calendar-effects:8}
Le lundi fait partie des nombreuses règles calendaires qu'on a examinées~; $t = 2.4$ ne résiste pas à une correction pour tests multiples. Avant de le négocier, il faut un mécanisme, une période
hors échantillon (out-of-sample) et les coûts de 52 allers-retours par an.
\end{solution}

\begin{solution}{pb:s1:seasonality-and-calendar-effects:1}
\begin{enumerate}
\item Un écart de rendement sur des jours définis par le calendrier~; rendements plus élevés les derniers et les premiers jours du mois~; récurrence des rendements relatifs à un même
moment de chaque année.
\item $+8.8~\%$ en janvier, $+7.3~\%$ en décembre, $-5.3~\%$ en mars et $-4.2~\%$ en avril par rapport à la tendance~: chauffage d'hiver, reconstitution des stocks au printemps.
\item Octobre et novembre sont les plus faibles ($t = -1.71$, $-2.17$), rien après correction (plus petite p-valeur ajustée 0.36).
\item Des rendements anormaux persistants autour des changements de semaine, de mois et d'année, ainsi qu'autour des jours fériés, sur quatre-vingt-dix ans.
\item Douze mois de 21 jours, neuf jours fériés, trente ans~; 8 points de base sur quatre jours de fin de mois, 15 sur les veilles de jours fériés~; une centaine de règles.
\item Selon leur écart attendu avec chaque effet injecté~: au moins 3 points de base suffisent à le faire ressortir.
\item Sans correction~: 19 règles de fin de mois, 0 de veille de jour férié, 8 fausses découvertes~; Bonferroni, Holm et Benjamini--Yekutieli~: 6, 0, 0~; Benjamini--Hochberg~: 11, 0, 3.
\item Neuf jours par an donnent 270 observations en trente ans~: un t attendu de 2.6, contre 1.5 dans cet échantillon.
\item Les rendements relatifs des actions se reproduisent au même mois du calendrier pour jusqu'à vingt décalages annuels.
\item Un rendement saisonnier mensuel de 1~\% par action~; le rendement moyen du même mois des cinq années précédentes.
\item Un IC de rang de 0.031 ($t = 3.8$)~; un spread entre déciles de 0.50~\% par mois, soit un ratio de Sharpe de 0.72 avant coûts.
\item Chaque action le montre une fois par an, face à un bruit mensuel de plusieurs pour cent.
\item \textbf{Résultat nommé.} L'\omterm{def:s1:seasonality-and-calendar-effects:tom}{effet de fin de mois} est retrouvé sous toutes les corrections (6 règles sous les corrections du taux d'erreur familial, 11 sous
Benjamini--Hochberg), l'effet de veille de jour férié sous aucune~; les fausses découvertes sont au nombre de 8 sans correction, 3 sous Benjamini--Hochberg et 0 sous Bonferroni,
Holm et Benjamini--Yekutieli.
\item Un mécanisme énoncé d'abord, un test qui compte chaque règle essayée, et assez d'observations.
\item Une correction du taux d'erreur familial (Holm) pour une recherche destinée à être négociée directement~; le taux de fausses découvertes pour un criblage suivi d'autres tests.
\item La fenêtre fixée à l'avance, chaque fenêtre essayée rapportée, les cours des contrats à terme et les coûts, une période hors échantillon (out-of-sample).
\item Que tester beaucoup de règles sur les mêmes données fait apparaître des effets qui n'existent pas.
\item L'effet de veille de jour férié.
\item Le livre~4 a établi les corrections~; ici, elles décident quels \omterm{def:s1:seasonality-and-calendar-effects:calendar}{effets calendaires} on retient.
\item Ceux qui disposent d'un mécanisme et qui résistent à une correction portant sur chaque règle essayée.
\end{enumerate}
\end{solution}

\begin{solution}{iq:s1:seasonality-and-calendar-effects:1}
Compter combien de règles calendaires ont été essayées (par moi et par la littérature), corriger pour elles, chercher un mécanisme, vérifier d'autres marchés et d'autres périodes,
et estimer ce qu'elle rapporterait après coûts~; le plus souvent, il faut la classer sans suite.
\end{solution}

\begin{solution}{iq:s1:seasonality-and-calendar-effects:2}
Le taux d'erreur familial (FWER) est la probabilité d'au moins une fausse découverte~; le taux de fausses découvertes (FDR) est la part attendue des découvertes qui sont
fausses. La première est plus stricte~; la seconde conserve davantage de vrais effets lorsqu'ils sont nombreux.
\end{solution}

\begin{solution}{iq:s1:seasonality-and-calendar-effects:3}
Les salaires, les cotisations de retraite et les afflux vers les fonds arrivent autour de la fin du mois et sont investis dans les jours qui suivent~; les \omterm{def:s1:index-rebalancing:fund}{fonds indiciels} se rééquilibrent en fin de mois. Vérifier les
données de flux par jour du mois, et si l'effet est plus marqué là où ces flux sont plus importants.
\end{solution}

\begin{solution}{iq:s1:seasonality-and-calendar-effects:4}
Les jours de négociation et les jours fériés de chaque bourse, tels qu'ils étaient chaque année~; les fermetures anticipées~; les échéances et les dates de règlement-livraison~; les fuseaux horaires~; et la correspondance entre les événements
connus à l'avance et ceux annoncés plus tard.
\end{solution}

\begin{solution}{iq:s1:seasonality-and-calendar-effects:5}
La saisonnalité au comptant est prévisible et intégrée dans la courbe des contrats à terme, si bien que les rendements de ces contrats n'ont pas à être saisonniers~; une \omterm{def:s1:seasonality-and-calendar-effects:calendar}{saisonnalité des rendements} des contrats à terme est
une prime de risque ou une mauvaise évaluation, et elle est bien plus difficile à trouver.
\end{solution}

\begin{solution}{iq:s1:seasonality-and-calendar-effects:6}
La différence des moyennes a pour erreur type $\sigma\sqrt{1/(fn) + 1/((1 - f)n)} = \sigma/\sqrt{f(1 - f)n}$, soit $E[t] = \delta\sqrt{f(1 - f)n}/\sigma$~; il
atteint $z$ lorsque $n = z^2\sigma^2/(\delta^2 f(1 - f))$.
\end{solution}
